% lark-formal.tex --- Lark (Lambda Affine Resource Kernel): Language Specification
%
% Part I: machine-checked STLC core (lcore proof files)
% Part II: full surface language (paper rules)
%
% Build:  tectonic lark-formal.tex
%
\documentclass[10pt,a4paper]{article}

\usepackage[T1]{fontenc}
\usepackage{lmodern}
\usepackage{amsmath,amssymb,amsthm}
\usepackage{mathpartir}
\usepackage[margin=2.2cm,top=2cm,bottom=2.5cm]{geometry}
\usepackage{microtype}
\usepackage{xcolor}
\usepackage{booktabs}
\usepackage[hidelinks]{hyperref}
\usepackage{alltt}

% -- notation

% Base types
\newcommand{\TInt}{\mathsf{Int}}
\newcommand{\TFlt}{\mathsf{Float}}
\newcommand{\TBool}{\mathsf{Bool}}
\newcommand{\TStr}{\mathsf{Str}}
\newcommand{\TUnit}{\mathsf{Unit}}
\newcommand{\TIO}{\mathsf{IO}}
\newcommand{\TCopy}{\mathsf{Copy}}

% Judgments
\newcommand{\hasty}[3]{#1 \vdash #2 : #3}       % Γ ⊢ e : τ
\newcommand{\inctx}[2]{#1 \ni #2}               % Γ ∋ τ
\newcommand{\isval}[2]{\mathsf{val}(#2)}         % val(e)  (τ suppressed in rules)
\newcommand{\IsVal}[2]{\mathsf{IsVal}(#1,\,#2)} % full form

% Reduction
\newcommand{\step}{\to}
\newcommand{\steps}{\to^{*}}
\newcommand{\eval}{\Downarrow}

% Substitution:  e[η]
% (η : SemEnv Γ maps each position to a closed term)
\newcommand{\subst}[2]{#1[#2]}

% Environments
\newcommand{\SemEnv}[1]{\mathsf{Env}(#1)}
\newcommand{\SemEnvD}[2]{\mathsf{Env}_{#1}(#2)}
\newcommand{\envnil}{\langle\,\rangle}
\newcommand{\envcons}[2]{(#1;\,#2)}

% Keywords
\newcommand{\klet}{\mathbf{let}}
\newcommand{\kin}{\mathbf{in}}
\newcommand{\kif}{\mathbf{if}}
\newcommand{\kthen}{\mathbf{then}}
\newcommand{\kelse}{\mathbf{else}}

% lcore file names
\newcommand{\lfile}[1]{\texttt{#1}}

% -- theorem environments
\theoremstyle{plain}
\newtheorem{theorem}{Theorem}
\theoremstyle{definition}
\newtheorem{definition}{Definition}
\theoremstyle{remark}
\newtheorem{remark}{Remark}

% -----------------------------------------------------------------------------
\title{\textbf{Lark: Language Specification}\\[6pt]
       \large Lambda Affine Resource Kernel\\[4pt]
       \normalsize Formal grammar, type rules, and operational semantics}
\author{}
\date{}

\begin{document}
\maketitle\thispagestyle{empty}

\noindent
This document is the formal specification of \textbf{Lark} (Lambda Affine
Resource Kernel): a small purely functional language with Hindley-Milner type
inference, affine ownership, traits, and a file-level module system.

\smallskip
\noindent\textbf{Part\,I} (\S\S\,\ref{sec:types}--\ref{sec:weaken}) is the
\emph{machine-checked} STLC core: types, typing rules, substitution, small-step
reduction, and a type-soundness theorem, encoded and verified in an MLTT kernel
with NbE.  Every formula has a counterpart in \texttt{proof/lark/}.

\smallskip
\noindent\textbf{Part\,II} (\S\S\,\ref{sec:types-ext}--\ref{sec:grammar})
extends to the complete language: algebraic data types, pattern matching, affine
ownership, traits, modules, and the full surface grammar.  These sections are
paper rules derived from the implementation; they are not machine-checked.

\smallskip
\noindent\textbf{Part\,III} (\S\S\,\ref{sec:int-defined}--\ref{sec:totality})
is an \emph{addendum}: the Phase\,8 rigor extensions (repository snapshot
\texttt{08/}).  Parts\,I--II specify the language as built through Phase\,7;
Part\,III records where Phase\,8 \emph{strengthens} or \emph{redefines} those
rules --- defined integer overflow, static exhaustiveness, and opt-in totality.
A reader following the build chronologically should read it as a delta,
exactly as the snapshots are.

\medskip\noindent\rule{\linewidth}{0.4pt}
{\centering\normalsize\bfseries Part\,I\quad Machine-Checked STLC Core\par}
\noindent\rule{\linewidth}{0.4pt}\medskip

% -----------------------------------------------------------------------------
\section{Types and Contexts}
\label{sec:types}

\paragraph{Types.}
\[
\tau,\sigma \;::=\; \TInt \;\mid\; \TFlt \;\mid\; \TBool \;\mid\; \TStr
            \;\mid\; \TUnit \;\mid\; \tau \to \sigma
\]

\paragraph{Contexts.}
A \emph{snoc-list} of types (de\,Bruijn style; innermost binding first):
\[
\Gamma \;::=\; \cdot \;\mid\; \Gamma,\tau
\]

\paragraph{Variable lookup.}
$\inctx{\Gamma}{\tau}$ witnesses that $\tau$ is bound at some position in $\Gamma$:
\begin{mathpar}
\inferrule{ }{\inctx{\Gamma{,}\tau}{\tau}}\;\textsc{Here}
\and
\inferrule{\inctx{\Gamma}{\tau}}{\inctx{\Gamma{,}\sigma}{\tau}}\;\textsc{There}
\end{mathpar}

% -----------------------------------------------------------------------------
\section{Typing Rules (\lfile{lark-typing.lcore})}
\label{sec:typing}

\noindent\textbf{Intrinsic encoding.}  A term of type
$\mathsf{Expr}(\Gamma,\tau)$ \emph{is} simultaneously an expression and a proof
that it has type~$\tau$ in context~$\Gamma$; there is no separate typing
relation.  The ten constructors correspond one-to-one with the rules below.

\begin{mathpar}
\inferrule{ }{\hasty{\Gamma}{n}{\TInt}}\;\textsc{LitInt}
\and
\inferrule{ }{\hasty{\Gamma}{f}{\TFlt}}\;\textsc{LitFlt}
\and
\inferrule{ }{\hasty{\Gamma}{b}{\TBool}}\;\textsc{LitBool}
\and
\inferrule{ }{\hasty{\Gamma}{s}{\TStr}}\;\textsc{LitStr}
\and
\inferrule{ }{\hasty{\Gamma}{()}{\TUnit}}\;\textsc{Unit}
\\
\inferrule{\inctx{\Gamma}{\tau}}
          {\hasty{\Gamma}{x}{\tau}}\;\textsc{Var}
\and
\inferrule{\hasty{\Gamma{,}\tau_1}{e}{\tau_2}}
          {\hasty{\Gamma}{\lambda x{:}\tau_1.\,e}{\tau_1\to\tau_2}}\;\textsc{Lam}
\\
\inferrule{\hasty{\Gamma}{f}{\tau_1\to\tau_2} \\
           \hasty{\Gamma}{a}{\tau_1}}
          {\hasty{\Gamma}{f\,a}{\tau_2}}\;\textsc{App}
\and
\inferrule{\hasty{\Gamma}{v}{\tau_1} \\
           \hasty{\Gamma{,}\tau_1}{b}{\tau_2}}
          {\hasty{\Gamma}{\klet\; x{=}v\;\kin\; b}{\tau_2}}\;\textsc{Let}
\and
\inferrule{\hasty{\Gamma}{c}{\TBool} \\
           \hasty{\Gamma}{t}{\tau} \\
           \hasty{\Gamma}{e}{\tau}}
          {\hasty{\Gamma}{\kif\;c\;\kthen\;t\;\kelse\;e}{\tau}}\;\textsc{If}
\end{mathpar}

% -----------------------------------------------------------------------------
\section{Semantic Environments and Substitution (\lfile{lark-subst.lcore})}
\label{sec:subst}

\paragraph{Semantic environments.}
$\SemEnv{\Gamma}$ maps each position in $\Gamma$ to a closed term of the
corresponding type:
\begin{align*}
\SemEnv{\cdot} &\;=\; \mathbf{1}\\
\SemEnv{\Gamma,\tau} &\;=\; \mathsf{Expr}(\cdot,\tau)\times\SemEnv{\Gamma}
\end{align*}
The definition is by recursion on $\Gamma$; the product unfolds \emph{definitionally},
which is the key to avoiding index-refinement in the kernel.

\paragraph{Variable lookup.}
$\mathsf{lookup} : \inctx{\Gamma}{\tau} \to \SemEnv{\Gamma} \to \mathsf{Expr}(\cdot,\tau)$
is defined by recursion on the derivation:
\[
  \mathsf{lookup}(\textsc{Here},\,(v;\eta)) = v
  \qquad
  \mathsf{lookup}(\textsc{There}(x),\,(v;\eta)) = \mathsf{lookup}(x,\,\eta)
\]

\paragraph{Open semantic environments.}
The closed-environment $\SemEnv{\Gamma}$ suffices for computing ground values
but not for substituting \emph{under a $\lambda$}.
We therefore generalise to an \emph{open} environment
\[
  \SemEnvD{\Delta}{\Gamma}
\]
which maps each position in $\Gamma$ to an \emph{open} term in context $\Delta$:
\begin{align*}
\SemEnvD{\Delta}{\cdot}           &\;=\; \mathbf{1}\\
\SemEnvD{\Delta}{\Gamma,\tau}     &\;=\; \mathsf{Expr}(\Delta,\tau)\times\SemEnvD{\Delta}{\Gamma}
\end{align*}
$\SemEnvD{\cdot}{\Gamma} = \SemEnv{\Gamma}$ definitionally.

Variable lookup in an open environment is identical in structure to the
closed case:
\[
  \mathsf{lookup}_\Delta(\textsc{Here},\,(v;\eta)) = v
  \qquad
  \mathsf{lookup}_\Delta(\textsc{There}(x),\,(v;\eta)) = \mathsf{lookup}_\Delta(x,\,\eta)
\]

\paragraph{Weakening an open environment.}
Given $\eta : \SemEnvD{\Delta}{\Gamma}$, inserting a new innermost binding
$\sigma$ into $\Delta$ gives $\overline{\eta}_\sigma : \SemEnvD{\Delta,\sigma}{\Gamma}$
by applying the kernel primitive \texttt{weaken} to each entry:
\[
  \overline{\eta}_\sigma \;=\; \mathsf{wk}_\sigma(\eta)
  \qquad\text{(see \S\ref{sec:weaken})}
\]

\paragraph{Semantic substitution.}
$\subst{e}{\eta}$ simultaneously replaces all free variables of $e$ with
the corresponding open terms in $\eta : \SemEnvD{\Delta}{\Gamma}$, producing
an open term in context $\Delta$ (\texttt{sub\_open}).
The ground case $\Delta = \cdot$ recovers closed substitution (\texttt{sub\_ground}).

\begin{center}
\renewcommand{\arraystretch}{1.35}
\begin{tabular}{@{}lll@{}}
\toprule
$e$ & $\subst{e}{\eta}$ & Note \\
\midrule
$n,\,f,\,b,\,s,\,()$ & $n,\,f,\,b,\,s,\,()$ & literals are closed \\
$x_{\textsc{here}}$ & $\mathsf{fst}\,\eta$ & innermost variable \\
$x_{\textsc{there}}$ & $\subst{x}{\mathsf{snd}\,\eta}$ & skip one binder \\
$f\,a$ & $\subst{f}{\eta}\;\subst{a}{\eta}$ & \\
$\klet\;x{=}v\;\kin\;b$
  & $\subst{b}{(\subst{v}{\eta};\,\eta)}$ & let-unrolling \\
$\kif\;c\;\kthen\;t\;\kelse\;e$
  & $\kif\;\subst{c}{\eta}\;\kthen\;\subst{t}{\eta}\;\kelse\;\subst{e}{\eta}$ & \\
$\lambda x{:}\tau_1.\,b$
  & $\lambda x{:}\tau_1.\,\subst{b}{(x_\mathsf{here};\,\overline{\eta}_{\tau_1})}$
  & extend \& weaken (see \S\ref{sec:weaken}) \\
\bottomrule
\end{tabular}
\end{center}

\paragraph{Value predicate.}
$\IsVal{\tau}{e}$ holds when the \emph{closed} expression $e : \mathsf{Expr}(\cdot,\tau)$
is already a value:
\begin{mathpar}
\inferrule{ }{\IsVal{\TInt}{n}}\;\textsc{ValInt}
\and
\inferrule{ }{\IsVal{\TFlt}{f}}\;\textsc{ValFlt}
\and
\inferrule{ }{\IsVal{\TBool}{b}}\;\textsc{ValBool}
\and
\inferrule{ }{\IsVal{\TStr}{s}}\;\textsc{ValStr}
\and
\inferrule{ }{\IsVal{\TUnit}{()}}\;\textsc{ValUnit}
\and
\inferrule{ }{\IsVal{\tau_1\to\tau_2}{\lambda x.\,e}}\;\textsc{ValLam}
\end{mathpar}

% -----------------------------------------------------------------------------
\section{Small-Step Reduction (\lfile{lark-step.lcore})}
\label{sec:step}

$e \step e'$ relates \emph{closed} terms of the same type $\tau$
(in the kernel: $\mathsf{Step}(\tau, e, e')$).

\medskip
\noindent\textbf{Base rules:}
\begin{mathpar}
\inferrule{ }
  {(\lambda x.\,b)\;v \;\step\; \subst{b}{\envcons{v}{\envnil}}}\;\textsc{$\beta$}
\and
\inferrule{ }
  {\kif\;\mathsf{true}\;\kthen\;t\;\kelse\;e \;\step\; t}\;\textsc{IfT}
\and
\inferrule{ }
  {\kif\;\mathsf{false}\;\kthen\;t\;\kelse\;e \;\step\; e}\;\textsc{IfF}
\and
\inferrule{\IsVal{\tau_1}{v}}
  {\klet\;x{=}v\;\kin\;b \;\step\; \subst{b}{\envcons{v}{\envnil}}}\;\textsc{LetBeta}
\end{mathpar}

\noindent\textbf{Congruence rules} (call-by-value order):
\begin{mathpar}
\inferrule{f \step f'}
  {f\,a \;\step\; f'\,a}\;\textsc{App1}
\and
\inferrule{\IsVal{\tau_1\to\tau_2}{f} \\ a \step a'}
  {f\,a \;\step\; f\,a'}\;\textsc{App2}
\and
\inferrule{c \step c'}
  {\kif\;c\;t\;e \;\step\; \kif\;c'\;t\;e}\;\textsc{IfCong}
\and
\inferrule{v \step v'}
  {\klet\;x{=}v\;\kin\;b \;\step\; \klet\;x{=}v'\;\kin\;b}\;\textsc{LetCong}
\end{mathpar}

\begin{remark}[Type preservation]
$\step$ is typed: $\mathsf{Step}(\tau,e,e')$ carries $\tau$ as an
explicit index, so both endpoints have the same type by construction.
Type preservation requires no separate proof.
\end{remark}

% -----------------------------------------------------------------------------
\section{Multi-Step Reduction and Evaluation (\lfile{lark-preservation})}
\label{sec:preservation}

\paragraph{Multi-step reduction.}
$e \steps e'$ is the reflexive-transitive closure of $\step$:
\begin{mathpar}
\inferrule{ }{e \steps e}\;\textsc{Refl}
\and
\inferrule{e \step e' \\ e' \steps e''}{e \steps e''}\;\textsc{Cons}
\end{mathpar}

\paragraph{Big-step evaluation.}
$e \eval v$ derives from $\step$ by requiring the endpoint to be a value:
\begin{mathpar}
\inferrule{\IsVal{\tau}{v}}{v \eval v}\;\textsc{Done}
\and
\inferrule{e \step e' \\ e' \eval v}{e \eval v}\;\textsc{Step}
\end{mathpar}

\paragraph{Main theorem.}

\begin{theorem}[Evaluation soundness]\label{thm:sound}
If\/ $e \eval v$ then $\IsVal{\tau}{v}$.
\end{theorem}
\begin{proof}
By induction on the derivation of $e \eval v$.
\begin{itemize}\itemsep2pt
  \item \textsc{Done}: $e = v$ and $\IsVal{\tau}{v}$ is a direct hypothesis.
  \item \textsc{Step}: $e \step e'$ and $e' \eval v$.
        By the induction hypothesis on $e' \eval v$, we obtain
        $\IsVal{\tau}{v}$. \qedhere
\end{itemize}
\end{proof}

\noindent
In the kernel, this is \texttt{eval\_produces\_val}, proved by
\texttt{indrec Eval} with motive $M(t,e,v,-) = \mathsf{IsVal}(t,v)$.
The \textsc{Done} case returns the \texttt{IsVal} witness directly;
the \textsc{Step} case passes through the induction hypothesis.
The normal form of the proof is:
\begin{center}
\small
$\texttt{indrec Eval}$\\[2pt]
$\quad(\lambda t\,e\,v\,\_.\;\mathsf{IsVal}(t,v))\;
(\lambda t\,v\,\mathit{iv}.\;\mathit{iv})\;
(\lambda t\,e\,e'\,v\,s\,\mathit{ev'}\,\mathit{ih}.\;\mathit{ih})$
\end{center}

% -----------------------------------------------------------------------------
\section{The \texttt{weaken} Primitive}
\label{sec:weaken}

The ELam case of \texttt{sub\_open} requires lifting each entry
$e_i : \mathsf{Expr}(\Delta,\tau_i)$ of $\eta : \SemEnvD{\Delta}{\Gamma}$
to $\mathsf{Expr}(\Delta,\sigma,\tau_i)$, so that the extended environment
has type $\SemEnvD{\Delta,\sigma}{\Gamma}$.  This operation is
\emph{weakening}:
\[
  \mathsf{weaken} : \forall\,\sigma\,\Delta\,\tau.\;
  \mathsf{Expr}(\Delta,\tau) \;\longrightarrow\; \mathsf{Expr}(\Delta,\sigma,\tau).
\]

\paragraph{Definitional obstacle.}
A definition of \texttt{weaken} by \texttt{indrec Expr} within the kernel
meets a definitional obstacle in the ELam case.  Given
$\mathsf{ELam}(\Delta,\tau_\lambda,\tau_2,\mathit{body})$ with
$\mathit{body}:\mathsf{Expr}(\mathsf{ext}(\tau_\lambda,\Delta),\tau_2)$,
the weakened lambda requires
$\mathit{body}':\mathsf{Expr}(\mathsf{ext}(\tau_\lambda,\mathsf{ext}(\sigma,\Delta)),\tau_2)$,
whereas a recursive call
$\mathsf{weaken}(\sigma,\mathsf{ext}(\tau_\lambda,\Delta),\tau_2,\mathit{body})$
produces type $\mathsf{Expr}(\mathsf{ext}(\sigma,\mathsf{ext}(\tau_\lambda,\Delta)),\tau_2)$.
Unifying these requires
\[
  \mathsf{ext}(\sigma,\,\mathsf{ext}(\tau_\lambda,\,\Delta))
  \;=\;
  \mathsf{ext}(\tau_\lambda,\,\mathsf{ext}(\sigma,\,\Delta)),
\]
which holds propositionally but not \emph{definitionally} in the kernel.

\paragraph{Kernel primitive.}
\texttt{weaken} is therefore a kernel primitive (\texttt{TM\_WEAKEN} in
\texttt{eval.c}), operating directly on the \texttt{VL\_INDCON} tree
representing an $\mathsf{Expr}$ value, via a de~Bruijn \emph{cutoff} $d$:
\begin{itemize}\itemsep2pt
\item \textbf{\texttt{weaken\_ctx}$(a,\Gamma,d)$} inserts type $a$ at
  depth $d$ in $\Gamma$; $d=0$ prepends, $d>0$ recurses under existing
  binders.
\item \textbf{\texttt{shift\_var\_val}$(v,a,d)$} shifts a $\mathsf{Var}$
  node; at $d=0$ a \textsc{Here} becomes \textsc{There}, at $d>0$ only
  the tail context is updated.
\item \textbf{\texttt{weaken\_expr\_val}$(e,a,d)$} recurses on $e$,
  updating every context-index argument and every variable reference; the
  ELam case increments $d$ to protect the lambda's own parameter.
\end{itemize}
All context-index arguments are updated consistently throughout the
constructed \texttt{VL\_INDCON} tree, so the result has the correct type
without appealing to any commutativity equation.

\paragraph{Syntax.} In \texttt{.lcore} files:
\begin{center}
\texttt{weaken}\ $\sigma$\ $\Delta$\ $\tau$\ $e$
$\;\;:\;\;$\texttt{Expr (ext}~$\sigma$~$\Delta$\texttt{)}~$\tau$
\end{center}

\paragraph{The ELam case of substitution.}
With \texttt{weaken} in place, the ELam row of the substitution table
(\S\ref{sec:subst}) is:
\begin{align*}
  \subst{\lambda x{:}\tau_1.\,b}{\eta}
  &\;=\;
  \lambda x{:}\tau_1.\;
  \subst{b}{\bigl(x_{\mathsf{here}};\;\mathsf{wk}_{\tau_1}(\eta)\bigr)}
\end{align*}
where $\mathsf{wk}_{\tau_1}(\eta)$ lifts each $e_i:\mathsf{Expr}(\Delta,\tau_i)$
in $\eta$ to $\mathsf{Expr}(\Delta,\tau_1,\tau_i)$ via \texttt{weaken}
(\texttt{weaken\_semctx\_d} in the source), and
$x_{\mathsf{here}}=\mathsf{EVar}(\Delta,\tau_1,\textsc{Here})$.
No commutativity equation is required.

As a concrete instance:
\begin{align*}
\texttt{sub\_ground}\ (\Gamma{=}\TInt)\ (\tau{=}\TBool{\to}\TBool)
\ (\lambda x{:}\TBool.\,x)\ [42]
\;=\;
\lambda x{:}\TBool.\,x
\end{align*}
(The lambda parameter remains at index~0 in the empty output context.)

\begin{remark}[Definitional equality of stuck \texttt{weaken} terms]
When the argument of \texttt{weaken} is neutral, the evaluator suspends
and records a \texttt{SP\_WEAKEN} spine entry.  Two terms
$\mathsf{weaken}(\sigma_1,\Delta,\tau,x)$ and
$\mathsf{weaken}(\sigma_2,\Delta,\tau,x)$ at the same neutral $x$ are
definitionally equal if and only if $\sigma_1$, $\Delta$, and $\tau$
each agree definitionally.  The conversion checker (\texttt{conv\_spine}
in \texttt{check.c}) compares all three parameters.
\end{remark}


\medskip\noindent\rule{\linewidth}{0.4pt}
{\centering\normalsize\bfseries Part\,II\quad Full Language Specification\par}
\noindent\rule{\linewidth}{0.4pt}\medskip

% -----------------------------------------------------------------------------
\section{Extended Types}
\label{sec:types-ext}

\paragraph{Types.}
The STLC core (\S\ref{sec:types}) uses six base types and function types.
The full language adds tuples, algebraic data types, type variables, and the
affine resource type $\TIO$:
\[
\tau,\sigma \;::=\;
  \TInt \mid \TFlt \mid \TBool \mid \TStr \mid \TUnit \mid \TIO
  \;\mid\; \tau\to\sigma
  \;\mid\; (\tau_1,\ldots,\tau_n)\;\;(n\ge 2)
  \;\mid\; \alpha
  \;\mid\; T \;\mid\; T(\tau_1,\ldots,\tau_n)
\]
$T$ ranges over declared type constructors; $\alpha$ over type variables
(lowercase names in source).

\paragraph{Type schemes.}
Let-generalisation (HM Algorithm W) produces type schemes
$\forall\bar\alpha.\,\tau$ at top-level \texttt{fn} and \texttt{let}
boundaries.  At each use site the bound variables are replaced by fresh
monotypes (instantiation).

\paragraph{Built-in types.}
\begin{center}
\renewcommand{\arraystretch}{1.2}
\begin{tabular}{@{}lll@{}}
\toprule
Type & Arity & $\TCopy$ \\
\midrule
$\TInt,\TFlt,\TBool,\TStr,\TUnit$       & 0 & yes \\
$\TIO$                                   & 0 & no (affine) \\
$\mathsf{List}(\alpha)$                  & 1 & iff $\TCopy(\alpha)$ \\
$\mathsf{Result}(\alpha,\beta)$          & 2 & iff $\TCopy(\alpha)\wedge\TCopy(\beta)$ \\
$(\tau_1,\ldots,\tau_n)$                 & $n\ge 2$ & iff $\forall i.\,\TCopy(\tau_i)$ \\
\bottomrule
\end{tabular}
\end{center}

% -----------------------------------------------------------------------------
\section{Affine Discipline}
\label{sec:affine}

Lark's ownership model is \emph{affine}: a value of a non-\texttt{Copy} type
may be used at most once.  Types opt into free duplication by satisfying the
$\TCopy$ predicate.

\paragraph{The \texttt{Copy} predicate.}
$\TCopy(\tau)$ is defined by structural recursion on $\tau$:
\begin{center}
\renewcommand{\arraystretch}{1.25}
\begin{tabular}{@{}lll@{}}
\toprule
$\tau$ & $\TCopy(\tau)$ & Note \\
\midrule
$\TInt,\TFlt,\TBool,\TStr,\TUnit$ & yes & built-in \\
$\tau\to\sigma$                    & yes & functions are code, not resources \\
$(\tau_1,\ldots,\tau_n)$           & $\TCopy(\tau_1)\wedge\cdots\wedge\TCopy(\tau_n)$ & \\
$\alpha$ (type variable)           & yes & conservative; prevents false positives \\
$T$ (nullary)                      & yes iff \texttt{impl Copy for T} declared & \\
$T(\tau_1,\ldots,\tau_n)$          & yes iff $T\in\mathsf{CopyTypes}\wedge\forall i.\,\TCopy(\tau_i)$ & \\
$\TIO$                             & no & no \texttt{Copy} impl \\
\bottomrule
\end{tabular}
\end{center}
\noindent $\mathsf{CopyTypes}$ is initialised to the built-in Copy types and
extended by each \texttt{impl Copy for T} declaration in the program.

\paragraph{Affine tracking.}
The type checker maintains a partial map
$\Omega : \mathsf{Var}\rightharpoonup\mathbb{N}$ (the \emph{affine
environment}) from variable names to use counts.

\begin{itemize}\itemsep3pt
\item \textbf{Bind.}  A parameter $x:\tau$ or let-binding $x=e:\tau$ with
  $\lnot\TCopy(\tau)$ is added to $\Omega$ with $\Omega(x)=0$.  Parameters
  with $\TCopy(\tau)$ are not tracked.
\item \textbf{Pattern binders.}  Variables introduced by match arms are not
  tracked; their types may be unresolved at check time and receive
  conservative Copy treatment.
\item \textbf{Globals.}  Top-level function names and built-in names are
  never tracked.
\item \textbf{Use.}  Each occurrence of tracked $x$ increments $\Omega(x)$.
  $\Omega(x)>1$ raises \textsc{AffineError}.  $\Omega(x)=0$ is permitted
  (affine values may be discarded).
\item \textbf{Shadow.}  A \texttt{let} rebinding a tracked name $x$ removes
  the old entry and creates a new one iff the new binding type is non-Copy.
  This allows the IO-threading idiom
  $\klet\;\mathsf{io}=\mathsf{print}(\mathsf{io},s)\;\kin\;\ldots$ without
  error.
\end{itemize}

% -----------------------------------------------------------------------------
\section{Extended Typing Rules}
\label{sec:typing-ext}

The judgment $\hasty{\Gamma}{e}{\tau}$ is the same as in
\S\ref{sec:typing}.  The affine constraint from \S\ref{sec:affine} applies
as a side condition on all rules: every tracked variable used in $e$ must
have use count $\le 1$ after checking.

\paragraph{Tuple introduction.}
\begin{mathpar}
\inferrule{\hasty{\Gamma}{e_1}{\tau_1} \\ \cdots \\ \hasty{\Gamma}{e_n}{\tau_n}}
          {\hasty{\Gamma}{(e_1,\ldots,e_n)}{(\tau_1,\ldots,\tau_n)}}\;\textsc{Tuple}
\quad(n\ge 2)
\end{mathpar}

\paragraph{Constructor introduction.}
Each ADT constructor $C$ is registered in the global constructor environment
$\Sigma$ with a type scheme.  Nullary constructors have scheme
$\forall\bar\alpha.\;T(\bar\alpha)$; constructors with payload have scheme
$\forall\bar\alpha.\;\tau_1\to\cdots\to\tau_k\to T(\bar\alpha)$:
\begin{mathpar}
\inferrule{C:\forall\bar\alpha.\;\tau_1\to\cdots\to\tau_k\to T(\bar\alpha)\;\in\;\Sigma}
          {\hasty{\Gamma}{C}{\tau_1[\bar\sigma/\bar\alpha]\to\cdots\to\tau_k[\bar\sigma/\bar\alpha]\to T(\bar\sigma)}}\;\textsc{Con}
\end{mathpar}
Constructor application is curried; \textsc{App} from \S\ref{sec:typing}
handles each argument.

\paragraph{Pattern matching.}
$\mathsf{binds}(p,\tau)$ is the environment of variable bindings introduced
by pattern $p$ at type $\tau$:
\begin{align*}
\mathsf{binds}(\_,\;\tau)
  &= \emptyset \\
\mathsf{binds}(x,\;\tau)
  &= \{x:\tau\} \\
\mathsf{binds}(n,\;\tau)
  &= \emptyset \quad(\text{literal}) \\
\mathsf{binds}(C(p_1,\ldots,p_k),\;T(\bar\sigma))
  &= {\textstyle\biguplus_i}\;\mathsf{binds}(p_i,\,\tau_i[\bar\sigma/\bar\alpha]) \\
\mathsf{binds}((p_1,\ldots,p_n),\;(\sigma_1,\ldots,\sigma_n))
  &= {\textstyle\biguplus_i}\;\mathsf{binds}(p_i,\,\sigma_i)
\end{align*}
where $\biguplus$ requires disjoint domains.
\begin{mathpar}
\inferrule{
  \hasty{\Gamma}{e}{\tau} \\\\
  \forall i:\;\hasty{\Gamma,\,\mathsf{binds}(p_i,\tau)}{e_i}{\sigma}
}{
  \hasty{\Gamma}{\mathtt{match}\;e\;\mathtt{with}\;\overline{|\;p_i\Rightarrow e_i}\;\mathtt{end}}{\sigma}
}\;\textsc{Match}
\end{mathpar}
All arms share the same result type $\sigma$ (unified by HM).  Arms are
tested in source order; a non-exhaustive match is a runtime error.

\paragraph{Function declarations and let-recursion.}
A \texttt{fn} declaration adds its own name to scope before checking the
body, enabling direct recursion:
\begin{mathpar}
\inferrule{
  \Gamma,\;f:\tau_1\to\cdots\to\tau_k\to\rho,\;
  x_1:\tau_1,\;\ldots,\;x_k:\tau_k \;\vdash\; e:\rho
}{
  \Gamma \;\vdash\;
  \mathtt{fn}\;f(x_1{:}\tau_1,\ldots,x_k{:}\tau_k){:}\rho=e
  \;\dashv\;
  f:\forall\bar\alpha.\,\tau_1\to\cdots\to\tau_k\to\rho
}\;\textsc{FnDecl}
\end{mathpar}
\emph{Mutual recursion}: all top-level \texttt{fn} declarations carrying
complete type annotations are added to $\Gamma$ simultaneously before any
body is checked (pre-registration pass).  Functions without full annotations
use the single-function let-rec path above.

\paragraph{Operators.}
\begin{mathpar}
\inferrule{\hasty{\Gamma}{e_1}{\alpha} \\ \hasty{\Gamma}{e_2}{\alpha} \\
           \alpha\in\{\TInt,\TFlt\}}
          {\hasty{\Gamma}{e_1\mathbin{\mathit{op}}e_2}{\alpha}}\;\textsc{Arith}
\quad(\mathit{op}\in\{+,-,*,/\})
\and
\inferrule{\hasty{\Gamma}{e_1}{\alpha} \\ \hasty{\Gamma}{e_2}{\alpha}}
          {\hasty{\Gamma}{e_1\mathbin{\mathit{cmp}}e_2}{\TBool}}\;\textsc{Cmp}
\quad(\mathit{cmp}\in\{{=}{=},{\ne},{<},{\le},{>},{\ge}\})
\and
\inferrule{\hasty{\Gamma}{e_1}{\TBool} \\ \hasty{\Gamma}{e_2}{\TBool}}
          {\hasty{\Gamma}{e_1\mathbin{\mathit{bop}}e_2}{\TBool}}\;\textsc{Bool}
\quad(\mathit{bop}\in\{\mathtt{and},\mathtt{or}\})
\end{mathpar}

\paragraph{IO operations.}
$\TIO$ is a pre-declared affine type with two built-in operations in $\Gamma$:
\[
\mathsf{print}:\TIO\to\TStr\to\TIO
\qquad
\mathsf{read}:\TIO\to(\TIO,\TStr)
\]
Since $\lnot\TCopy(\TIO)$, the affine discipline ensures each IO token
threads linearly through the call chain.

% -----------------------------------------------------------------------------
\section{Extended Operational Semantics}
\label{sec:step-ext}

The small-step relation $e\step e'$ and value predicate $\IsVal{\tau}{v}$
from \S\ref{sec:step} are extended with tuples, constructor values, and
pattern-match reduction.

\paragraph{Extended value forms.}
\begin{mathpar}
\inferrule{\IsVal{\tau_i}{v_i}\;\text{for each }i}
          {\IsVal{(\tau_1,\ldots,\tau_n)}{(v_1,\ldots,v_n)}}\;\textsc{ValTup}
\and
\inferrule{C:\forall\bar\alpha.\;\tau_1\to\cdots\to\tau_k\to T(\bar\alpha)\;\in\;\Sigma \\
           \IsVal{\tau_i[\bar\sigma/\bar\alpha]}{v_i}\;\text{for each }i}
          {\IsVal{T(\bar\sigma)}{C\,v_1\cdots v_k}}\;\textsc{ValCon}
\end{mathpar}

\paragraph{Pattern-match function.}
$\mathsf{match}(v,p)$ returns a substitution $\theta$ or $\mathsf{fail}$:
\begin{align*}
\mathsf{match}(v,\,\_)
  &= \{\} \\
\mathsf{match}(v,\,x)
  &= \{x\mapsto v\} \\
\mathsf{match}(n,\,n)
  &= \{\} \quad(\text{same literal}) \\
\mathsf{match}(C\,v_1\cdots v_k,\;C(p_1,\ldots,p_k))
  &= {\textstyle\biguplus_i}\;\mathsf{match}(v_i,p_i) \\
\mathsf{match}(v,\;C'(\ldots))
  &= \mathsf{fail}\quad(C'\ne C\text{, or arity mismatch}) \\
\mathsf{match}((v_1,\ldots,v_n),\;(p_1,\ldots,p_n))
  &= {\textstyle\biguplus_i}\;\mathsf{match}(v_i,p_i)
\end{align*}

\paragraph{Match reduction.}
\begin{mathpar}
\inferrule{
  \mathsf{match}(v,\,p_j)=\theta \\
  \forall i<j:\;\mathsf{match}(v,\,p_i)=\mathsf{fail}
}{
  \mathtt{match}\;v\;\mathtt{with}\;\overline{|\;p_i\Rightarrow e_i}\;\mathtt{end}
  \;\step\;e_j[\theta]
}\;\textsc{MatchArm}
\and
\inferrule{e\step e'}
          {\mathtt{match}\;e\;\mathtt{with}\;\ldots\;\mathtt{end}
           \;\step\;
           \mathtt{match}\;e'\;\mathtt{with}\;\ldots\;\mathtt{end}}\;\textsc{MatchCong}
\end{mathpar}

\paragraph{Tuple congruence.}
\begin{mathpar}
\inferrule{e_i\step e_i' \\ \IsVal{\tau_j}{v_j}\;\text{for }j<i}
          {(v_1,\ldots,v_{i-1},e_i,e_{i+1},\ldots)
           \;\step\;
           (v_1,\ldots,v_{i-1},e_i',e_{i+1},\ldots)}\;\textsc{TupCong}
\end{mathpar}
Constructor arguments reduce left-to-right via the existing
\textsc{App1}/\textsc{App2} rules (\S\ref{sec:step}); constructor
application is curried.

% -----------------------------------------------------------------------------
\section{Trait System}
\label{sec:traits}

Traits provide bounded polymorphism.  A trait $T[\alpha]$ declares a family
of methods; an \texttt{impl} provides those methods for a concrete type.

\paragraph{Trait declaration.}
\[
\mathtt{trait}\;T\;\alpha = \bigl\{\;\mathtt{fn}\;m_1:\tau_1[\alpha]\;;\;\cdots\;;\;\mathtt{fn}\;m_k:\tau_k[\alpha]\;\bigr\}
\]
registers $m_1,\ldots,m_k$ in the global trait environment $\Phi$.

\paragraph{Implementation.}
\[
\mathtt{impl}\;T\;\tau_0\;\mathtt{for}\;T_0 =
\bigl\{\;\mathtt{fn}\;m_i(\ldots)=e_i\;\bigr\}
\]
provides each method $m_i$ for the concrete type $T_0$.  After
type-checking, $(m_i,T_0)$ is entered in the dispatch table
$\mathcal{D}:(\mathsf{Name}\times\mathsf{Type})\to\mathsf{Closure}$.

\paragraph{Bounded functions.}
A function may impose trait constraints on type variables:
\[
\mathtt{fn}\;f\;[T_1\,\alpha_1,\ldots,T_n\,\alpha_n]
  (x_1:\sigma_1,\ldots):\rho = e
\]
The body is checked in $\Gamma$ extended with $\alpha_i\in T_i$ for each
bound.  Within the body, any method $m$ of $T_i$ may be applied to a
term of type $\alpha_i$.

\paragraph{Method dispatch.}
At call time the concrete type $T_0$ of the receiver is determined and
$\mathcal{D}(m,T_0)$ is looked up.  The compiler generates one dispatch stub
per trait method; the stub inspects the constructor tag of the argument
and branches to the appropriate implementation.

\begin{remark}
Dispatch stubs operate on ADT constructor tags.  Primitive types
($\TInt$, $\TFlt$, $\TBool$, $\TStr$) carry no heap tag; trait dispatch
for these types is routed through named runtime stubs rather than the
tag-checking path.
\end{remark}

% -----------------------------------------------------------------------------
\section{Module System}
\label{sec:modules}

Each \texttt{.lark} source file defines exactly one module.

\paragraph{Structure.}
A file begins with its module declaration, followed by any import
declarations, then top-level declarations:
\begin{align*}
&\mathtt{module}\;M \\
&\mathtt{import}\;N\;\mathtt{exposing}\;(x_1,\ldots,x_k) \\
&\textit{decl}_1\;\;\cdots\;\;\textit{decl}_n
\end{align*}
Top-level declarations may be prefixed with \texttt{export} to make them
visible to importing modules.

\paragraph{Scoping.}
The type checker processes modules in dependency order.  Importing module $N$
from module $M$ prepends the exported declarations of $N$ to $M$'s
declaration list before type-checking.  Imported function bodies are
therefore visible to all subsequent passes: closure conversion, TAC lowering,
and assembly.

\paragraph{Restrictions.}
Circular imports are not supported.  Every name used in a declaration must
be declared in the same file or imported explicitly.

% -----------------------------------------------------------------------------
\section{Surface Grammar}
\label{sec:grammar}

The complete Lark grammar in ISO EBNF\@.  \texttt{name} begins with a
lowercase letter; \texttt{Name} with uppercase; \texttt{\_} is a reserved
wildcard token.  In EBNF, \texttt{\{}\,$x$\,\texttt{\}} denotes zero or more
repetitions of $x$; \texttt{[}\,$x$\,\texttt{]} denotes an optional $x$;
quoted strings are literal tokens.

\medskip\noindent\textbf{Program structure}
{\small\begin{alltt}
program     = module_decl \{ import_decl \} \{ top_decl \} ;
module_decl = "module" Name ;
import_decl = "import" Name [ "exposing" "(" name_list ")" ] ;
name_list   = export_name \{ "," export_name \} ;
export_name = name | Name ;
top_decl    = [ "export" ] decl ;
decl        = fn_decl | let_decl | type_decl | trait_decl | impl_decl ;
\end{alltt}}

\medskip\noindent\textbf{Declarations}
{\small\begin{alltt}
fn_decl      = "fn" name [ "[" bounds "]" ] "(" [ params ] ")"
               [ ":" type ] "=" expr ;
let_decl     = "let" name [ ":" type ] "=" expr ;
type_decl    = "type" Name \{ name \} "=" type_body ;
trait_decl   = "trait" Name \{ name \} "=" "\{" \{ trait_method \} "\}" ;
impl_decl    = "impl" Name \{ atom_type \} "for" atom_type
               "=" "\{" \{ impl_method \} "\}" ;
params       = param \{ "," param \} ;  param = name [ ":" type ] ;
bounds       = bound \{ "," bound \} ;  bound = Name name ;
trait_method = "fn" name ":" type ;
impl_method  = "fn" name "(" [ params ] ")" "=" expr ;
\end{alltt}}

\medskip\noindent\textbf{Expressions}
{\small\begin{alltt}
expr        = let_expr | if_expr | match_expr | lambda_expr | infix_expr ;
let_expr    = "let" name [ ":" type ] "=" expr "in" expr ;
if_expr     = "if" expr "then" expr "else" expr ;
match_expr  = "match" expr "with" \{ "|" pattern "=>" expr \} "end" ;
lambda_expr = "fn" "(" [ params ] ")" "=>" expr ;
infix_expr  = unary_expr \{ op unary_expr \} ;
unary_expr  = unary_op unary_expr | apply_expr ;
unary_op    = "-" | "not" ;
apply_expr  = atom \{ "(" [ args ] ")" \} ;
args        = expr \{ "," expr \} ;
atom        = literal | name | Name | "(" expr ")"
            | "(" expr "," expr \{ "," expr \} ")" ;
\end{alltt}}

\medskip\noindent\textbf{Types and patterns}
{\small\begin{alltt}
type      = fn_type ;
fn_type   = atom_type "->" fn_type | atom_type ;
atom_type = Name | Name "(" type \{ "," type \} ")" | name | "()"
          | "(" type ")" | "(" type "," type \{ "," type \} ")" ;
type_body = type | "|" Name [ "of" type ] \{ "|" Name [ "of" type ] \} ;

pattern   = "_" | literal | name
          | Name [ "(" pattern \{ "," pattern \} ")" ]
          | "(" pattern "," pattern \{ "," pattern \} ")" ;
\end{alltt}}

\medskip\noindent\textbf{Operators and literals}
{\small\begin{alltt}
op      = "+" | "-" | "*" | "/" | "==" | "!="
        | "<" | "<=" | ">" | ">=" | "and" | "or" ;
literal = integer | float | string | "true" | "false" | "()" ;
integer = digit \{ digit \} ;
float   = digit \{ digit \} "." digit \{ digit \} ;
string  = '"' \{ str_char \} '"' ;
\end{alltt}}

\medskip\noindent\textbf{Lexical}
{\small\begin{alltt}
name = lowercase \{ letter | digit | "_" \} ;
Name = uppercase \{ letter | digit | "_" \} ;
\end{alltt}}
\noindent Keywords (not valid as \texttt{name} or \texttt{Name}):
\texttt{and else end export false fn if impl import in let match
module not of or then trait true type where with}.

\bigskip
\noindent\rule{\linewidth}{0.4pt}
{\centering\normalsize\bfseries Part\,III\quad Addendum: Phase\,8 Rigor Extensions\par}
\noindent\rule{\linewidth}{0.4pt}\medskip

\noindent
This part specifies the Phase\,8 extensions (snapshot \texttt{08/}).  Like
Part\,II these are paper rules; unlike Part\,II each rule here is enforced by
an executable check in \texttt{08/src/} and exercised by the differential
test suite across all four backends.  Machine-checking them in the MLTT
kernel is future work.

% -----------------------------------------------------------------------------
\section{Defined Integer Semantics}
\label{sec:int-defined}

Through Phase\,7 the behaviour of $\TInt$ arithmetic outside
$[-2^{31},\,2^{31})$ was backend-dependent (the Python evaluators used
unbounded integers; the RV32 backend wrapped).  Phase\,8 \emph{defines} it:

\begin{definition}[Int values]
$\TInt$ values are the two's-complement 32-bit integers
$\mathbb{Z}_{32}=\{-2^{31},\ldots,2^{31}-1\}$.  The wrap function
$\mathsf{w}:\mathbb{Z}\to\mathbb{Z}_{32}$ is
\[
\mathsf{w}(n) \;=\; \bigl((n + 2^{31}) \bmod 2^{32}\bigr) - 2^{31}.
\]
\end{definition}

\paragraph{Arithmetic.}
For $n_1,n_2\in\mathbb{Z}_{32}$ and $\oplus\in\{+,-,\times\}$, the reduction
rule for primitive application computes in $\mathbb{Z}$ and wraps:
\[
n_1 \oplus n_2 \;\step\; \mathsf{w}(n_1 \oplus_{\mathbb{Z}} n_2).
\]

\paragraph{Division and remainder.}
Division truncates toward zero and follows the RISC-V M-extension
conventions, so no arithmetic operation is a stuck term:
\begin{align*}
n_1 \,/\, n_2 &\;\step\;
  \begin{cases}
    -1                & n_2 = 0\\
    -2^{31}           & n_1 = -2^{31},\; n_2 = -1\\
    \mathsf{trunc}(n_1/n_2) & \text{otherwise}
  \end{cases}
&
n_1 \,\%\, n_2 &\;\step\;
  \begin{cases}
    n_1               & n_2 = 0\\
    0                 & n_1 = -2^{31},\; n_2 = -1\\
    n_1 - n_2\cdot\mathsf{trunc}(n_1/n_2) & \text{otherwise}
  \end{cases}
\end{align*}

\begin{remark}
Wrap was chosen over trap because the compilation target decides: RV32
hardware wraps for free, so wrapping is the only semantics every backend can
implement at zero cost.  The consequence is visible in
\S\ref{sec:totality}: integer descent is not well-founded under wrapping,
and the totality judgment must refuse it.
\end{remark}

% -----------------------------------------------------------------------------
\section{Exhaustiveness}
\label{sec:exhaust}

Phase\,8 strengthens the \textsc{Match} typing rule of \S\ref{sec:typing-ext}
with a coverage side condition, checked by Maranget's usefulness algorithm
(\texttt{08/src/exhaust.py}).

\begin{definition}[Pattern matrix, specialisation, default]
Let $P$ be an $m\times n$ matrix of patterns.  For a constructor $c$ of
arity $k$, the \emph{specialisation} $\mathcal{S}(c,P)$ keeps each row whose
first pattern is $c(p_1,\ldots,p_k)$ (replacing it by
$p_1,\ldots,p_k$) or a wildcard/variable (replacing it by $k$ wildcards),
and drops the rest.  The \emph{default} $\mathcal{D}(P)$ keeps the rows
whose first pattern is a wildcard/variable, minus their first column.
\end{definition}

\begin{definition}[Usefulness, coverage]
$\mathcal{U}(P,\vec q)$ holds iff some value vector matches $\vec q$ but no
row of $P$.  A match on scrutinee type $\tau$ with arm patterns
$p_1,\ldots,p_m$ is \emph{exhaustive} iff
$\neg\,\mathcal{U}\bigl((p_1)\cdots(p_m)^{T},\;(\_)\bigr)$,
decided by structural recursion on $\mathcal{S}$ and $\mathcal{D}$,
driven by the column type's \emph{complete signature}: the declared
constructors for an ADT, $\{(),\}$ for $\TUnit$,
$\{\mathtt{true},\mathtt{false}\}$ for $\TBool$, the single $n$-tuple
constructor for products, and \emph{no} finite signature for $\TInt$,
$\TFlt$, $\TStr$, $\TIO$, functions, and type variables (these require a
wildcard or variable arm).
\end{definition}

\paragraph{Strengthened rule.}
\begin{mathpar}
\inferrule{
  \hasty{\Gamma}{e}{\tau} \\
  \hasty{\Gamma,\;\mathsf{binds}(p_i:\tau)}{e_i}{\sigma}\;\text{for each }i \\
  \mathsf{exhaustive}(\tau;\;p_1,\ldots,p_m)
}{
  \hasty{\Gamma}{\mathtt{match}\;e\;\mathtt{with}\;\overline{|\;p_i\Rightarrow e_i}\;\mathtt{end}}{\sigma}
}\;\textsc{Match$^{+}$}
\end{mathpar}

\begin{theorem}[Match progress, paper]
If a match expression is well typed under \textsc{Match$^{+}$} and its
scrutinee is a value $v$, then $\mathsf{match}(v,p_j)$ succeeds for some
arm $j$: the \textsc{MatchArm} rule of \S\ref{sec:step-ext} always applies,
and the match-failure trap emitted by the compiler is unreachable.
\end{theorem}

\begin{remark}
When the check fails, the usefulness recursion has \emph{constructed} a
value vector no arm matches; it is reported as a witness pattern
(e.g.\ \texttt{Cons(Err(\_), Nil)}), turning the rejection into a
diagnosis.
\end{remark}

% -----------------------------------------------------------------------------
\section{Totality}
\label{sec:totality}

A function may opt into a termination check:
\texttt{fn total} $f(x_1{:}\tau_1,\ldots,x_n{:}\tau_n) = e$.
The modifier is contextual --- \texttt{total} remains a valid identifier.

\begin{definition}[Structural order]
Write $y \prec x$ when $y$ is bound at depth ${\ge}1$ inside a constructor
or tuple pattern matched against $x$, or against some $z \prec x$.  The
relation is well-founded: every Lark value is a finite tree.
\end{definition}

\begin{definition}[Totality judgment]
$\mathsf{total}(f)$ holds iff
\begin{enumerate}
\item some fixed argument position $i$ decreases at \emph{every} recursive
      call $f(a_1,\ldots,a_n)$ in $e$: the argument $a_i$ is a variable
      with $a_i \prec x_i$;
\item every other name referenced in $e$ is known to terminate: a builtin,
      a constructor, a parameter of $f$ (totality is \emph{relative} --- the
      caller supplies total arguments), another \texttt{total} function, or
      a function whose call graph reaches no cycle;
\item $f$ itself appears only in call position, and $f$ is not part of a
      mutual-recursion cycle (a v1 restriction; lexicographic and
      size-change measures are future work).
\end{enumerate}
\end{definition}

\begin{theorem}[Relative termination, paper]
If $\mathsf{total}(f)$ and $f$ is applied to values whose function-typed
components are themselves terminating, evaluation of the application
terminates.
\end{theorem}

\begin{remark}
$f(n-1)$ on $\TInt$ is \emph{not} accepted as decreasing, and this is
soundness rather than incompleteness: under the wrapping semantics of
\S\ref{sec:int-defined}, descent from a negative $n$ never reaches a base
case at $0$ --- it wraps through $-2^{31}$.  The two Phase\,8 semantic
decisions are mutually constraining: defined overflow is what makes
integer descent genuinely ill-founded, and structural descent on data is
the discipline that remains.
\end{remark}

\end{document}
